\maketitle%与上面的\title对应
\begin{abstract}%摘要+摘要+摘要+摘要+摘要

	农作物种植结构优化是保障粮食安全与农户收益的关键问题。本文基于某乡村54块地块、41种作物的种植与销售数据，围绕"确定性—随机—相关随机"三层递进的建模框架展开系统研究，给出2024--2030年可直接落地的种植方案。

	\textbf{对于问题一}，将"地块$\times$季次$\times$作物"的安排抽象为0-1整数线性规划，完整刻画土地容量、水浇地二季复种、重茬回避、豆类三年轮作与销售上限等约束，分别求解"超产滞销"（情形1）与"超产按50\%折价出售"（情形2）两种销售规则下7年总净收益最大的方案。结果表明，情形1下2024--2030年总净收益约$\mathbf{3988.2}$万元，种植结构以粮食为主；情形2允许半价消化超产，模型大幅扩种高收益蔬菜，总净收益升至$\mathbf{5918.0}$万元，较情形1提高约$48.4\%$。

	\textbf{对于问题二}，针对销售量、亩产量、价格、成本的不确定性，按题目给定波动区间抽样$N=30$个等概率情景，建立以种植方案为第一阶段决策、情景销售为第二阶段决策的两阶段随机规划，并以"期望收益$-\lambda\cdot\mathrm{CVaR}_{0.90}$(损失)"为目标兼顾期望与尾部风险。结果表明，独立情景下期望收益约$\mathbf{4022.3}$万元，收益标准差仅$1.40\%$、最差情景偏离期望$3.69\%$，组合风险很低；不同风险系数$\lambda$下最优方案几乎不变，说明多样化已充分分散可对冲的个体风险。

	\textbf{对于问题三}，从天气共同因子、蔬菜市场因子、替代与互补消费、产销联动等角度构造170维作物冲击的相关矩阵，经Higham投影与Cholesky分解生成相关情景后重新优化。结果表明，期望收益与问题2基本持平（$\mathbf{4019.3}$万元），但收益标准差由$1.40\%$升至$2.46\%$、CVaR短缺口由$104.8$万元升至$173.8$万元，组合风险约放大$1.7$倍（跨种子复核$1.3\sim1.8$倍）；两种口径下最优种植方案几乎一致，说明相关结构主要影响风险测度而非方案本身，实际生产中须正视独立假设下被低估的尾部风险。

	本文建立了可复现、可检验的种植决策优化框架，各问方案均通过全部约束的程序化校验，可为该乡村及同类地区的种植规划提供决策参考。

	\keywords{0-1整数规划；两阶段随机规划；CVaR；作物种植优化；相关情景模拟}
\end{abstract}
